% --- Archon physics formalization source begin ---
% archon:physics
% archon:covers PhyXMiniProblems/problem_phyx_mini_0347.lean
% archon:source-report reports/phyx_mini/problem_phyx_mini_0347.source.json

\chapter{Physics problem phyx\_mini\_0347}
\label{ch:PhyXMiniProblems_problem_phyx_mini_0347}

\paragraph{Problem source.}
\#\# Physical scenario

The \textbackslash{}( pV \textbackslash{})-diagram in \textbackslash{}textbf\{figure\} shows a cycle of a heat engine that uses 0.250 mol of an ideal gas with \textbackslash{}( \textbackslash{}gamma = 1.40 \textbackslash{}). Process \textbackslash{}( ab \textbackslash{}) is adiabatic. (a) Find the pressure of the gas at point \textbackslash{}( a \textbackslash{}).

\#\# Question

What is the thermal efficiency of the engine?

\#\# Answer choices

- A: A: 30\textbackslash{}\%

- B: B: 31.9\textbackslash{}\%

- C: C: 35.9\textbackslash{}\%

- D: D: 32.3\textbackslash{}\%

\#\# Figure caption (auxiliary; use the image as primary evidence)

The image is a pressure-volume (p-V) diagram depicting a thermodynamic cycle. It has the following features:

- **Axes**: The vertical axis represents pressure (p) in atmospheres (atm), and the horizontal axis represents volume (V) in cubic meters (m³).

- **Points**: There are three labeled points on the graph: \textbackslash{}(a\textbackslash{}), \textbackslash{}(b\textbackslash{}), and \textbackslash{}(c\textbackslash{}).
  - Point \textbackslash{}(a\textbackslash{}) is at the top-left.
  - Point \textbackslash{}(b\textbackslash{}) is at the bottom-right.
  - Point \textbackslash{}(c\textbackslash{}) is at the bottom-left.

- **Lines and Arrows**:
  - A curved line connects points \textbackslash{}(a\textbackslash{}) and \textbackslash{}(b\textbackslash{}), with an arrow pointing from \textbackslash{}(a\textbackslash{}) to \textbackslash{}(b\textbackslash{}).
  - A horizontal line connects points \textbackslash{}(b\textbackslash{}) and \textbackslash{}(c\textbackslash{}), with an arrow pointing from \textbackslash{}(b\textbackslash{}) to \textbackslash{}(c\textbackslash{}).
  - A vertical line connects points \textbackslash{}(c\textbackslash{}) and \textbackslash{}(a\textbackslash{}), with an arrow pointing from \textbackslash{}(c\textbackslash{}) to \textbackslash{}(a\textbackslash{}).

- **Dotted Lines**: Vertical dotted lines extend from points \textbackslash{}(b\textbackslash{}) and \textbackslash{}(c\textbackslash{}) to the horizontal axis.

- **Labels**:
  - The point \textbackslash{}(a\textbackslash{}) is at the left, where there is a higher pressure that decreases in the direction toward \textbackslash{}(b\textbackslash{}).
  - The point \textbackslash{}(b\textbackslash{}) has a larger volume compared to point \textbackslash{}(c\textbackslash{}).

- **Numerical Values**: The pressure value labeled on the horizontal dashed line is 1.5 atm. The volume values are 0.0020 m³ and 0.0090 m³ at points \textbackslash{}(c\textbackslash{}) and \textbackslash{}(b\textbackslash{}), respectively.

The diagram illustrates a cycle involving compression and expansion, common in thermodynamics studies.

\#\# Dataset metadata

Laws of Thermodynamics; Physical Model Grounding Reasoning, Multi-Formula Reasoning

\paragraph{Recorded answer/context.}
B: B: 31.9\textbackslash{}\%

\paragraph{Figure/image path.}
/root/proposal\_for\_physic/science-mango/phyx\_mini\_run/phyx\_data/test\_image/347.png

\paragraph{Formalization target.}
create a compiling Lean file with sorry bodies at `PhyXMiniProblems/problem\_phyx\_mini\_0347.lean`. The Lean declarations must preserve the physical quantities, dimensions or dimensional roles, figure labels, governing-law hypotheses, and final relation expressed by this problem.
Use Mathlib/Physlib names found through LeanExplore where available. If a physics API is missing, introduce faithful local abstractions rather than scalar placeholder aliases.

\begin{theorem}[Physics formalization target]
\label{thm:physics:phyx_mini_0347:target}
\lean{PhyXMiniProblems.ProblemPhyXMini0347.thermalEfficiency_choiceB}
\uses{def:physics:phyx-mini-0347:phyxminiproblems-problemphyxmini0347-idealgasheatenginecycle, def:physics:phyx-mini-0347:phyxminiproblems-problemphyxmini0347-hasphysicalheatengineparameters, def:physics:phyx-mini-0347:phyxminiproblems-problemphyxmini0347-satisfiesidealgasheatenginelaws, def:physics:phyx-mini-0347:phyxminiproblems-problemphyxmini0347-hasheatenginefiguredata, def:physics:phyx-mini-0347:phyxminiproblems-problemphyxmini0347-isclosestefficiencychoice}
The assigned autoformalize agent should translate this physics problem into Lean declarations in the covered file, with theorem and lemma proof bodies written as `by sorry`.
\end{theorem}
\begin{proof}
This is an autoformalization task, not a proof task. Produce statements that can later be proved by the physics prover without weakening the physical model.
\end{proof}

% --- Archon declaration topology begin ---
\section{Declaration topology}
\begin{definition}[Gas Volume]
  \label{def:physics:phyx-mini-0347:phyxminiproblems-problemphyxmini0347-gasvolume}
  \lean{PhyXMiniProblems.ProblemPhyXMini0347.GasVolume}
  A physical gas volume, carrying the length-cubed dimension
\end{definition}
\begin{proof}
  This is the declared model definition of Gas Volume.
\end{proof}
\begin{definition}[pressure In Pascals]
  \label{def:physics:phyx-mini-0347:phyxminiproblems-problemphyxmini0347-pressureinpascals}
  \lean{PhyXMiniProblems.ProblemPhyXMini0347.pressureInPascals}
  The coherent-SI pressure readout, in pascals
\end{definition}
\begin{proof}
  This is the declared model definition of pressure In Pascals.
\end{proof}
\begin{definition}[pressure In Atmospheres]
  \label{def:physics:phyx-mini-0347:phyxminiproblems-problemphyxmini0347-pressureinatmospheres}
  \lean{PhyXMiniProblems.ProblemPhyXMini0347.pressureInAtmospheres}
  \uses{def:physics:phyx-mini-0347:phyxminiproblems-problemphyxmini0347-pressureinpascals}
  The pressure readout as a dimensionless multiple of one standard atmosphere
\end{definition}
\begin{proof}
  This is the declared model definition of pressure In Atmospheres.
\end{proof}
\begin{definition}[volume In Cubic Meters]
  \label{def:physics:phyx-mini-0347:phyxminiproblems-problemphyxmini0347-volumeincubicmeters}
  \lean{PhyXMiniProblems.ProblemPhyXMini0347.volumeInCubicMeters}
  \uses{def:physics:phyx-mini-0347:phyxminiproblems-problemphyxmini0347-gasvolume}
  The coherent-SI volume readout, in cubic metres
\end{definition}
\begin{proof}
  This is the declared model definition of volume In Cubic Meters.
\end{proof}
\begin{definition}[energy In Joules]
  \label{def:physics:phyx-mini-0347:phyxminiproblems-problemphyxmini0347-energyinjoules}
  \lean{PhyXMiniProblems.ProblemPhyXMini0347.energyInJoules}
  The coherent-SI energy readout, in joules
\end{definition}
\begin{proof}
  This is the declared model definition of energy In Joules.
\end{proof}
\begin{definition}[temperature In Kelvin]
  \label{def:physics:phyx-mini-0347:phyxminiproblems-problemphyxmini0347-temperatureinkelvin}
  \lean{PhyXMiniProblems.ProblemPhyXMini0347.temperatureInKelvin}
  The cycle uses the underlying absolute-temperature readout in kelvin
\end{definition}
\begin{proof}
  This is the declared model definition of temperature In Kelvin.
\end{proof}
\begin{definition}[State Label]
  \label{def:physics:phyx-mini-0347:phyxminiproblems-problemphyxmini0347-statelabel}
  \lean{PhyXMiniProblems.ProblemPhyXMini0347.StateLabel}
  The three equilibrium states labelled in the pressure-volume diagram
\end{definition}
\begin{proof}
  This is the declared model definition of State Label.
\end{proof}
\begin{definition}[Process Kind]
  \label{def:physics:phyx-mini-0347:phyxminiproblems-problemphyxmini0347-processkind}
  \lean{PhyXMiniProblems.ProblemPhyXMini0347.ProcessKind}
  The process types explicitly specified by the text and the diagram
\end{definition}
\begin{proof}
  This is the declared model definition of Process Kind.
\end{proof}
\begin{definition}[Thermodynamic State]
  \label{def:physics:phyx-mini-0347:phyxminiproblems-problemphyxmini0347-thermodynamicstate}
  \lean{PhyXMiniProblems.ProblemPhyXMini0347.ThermodynamicState}
  \uses{def:physics:phyx-mini-0347:phyxminiproblems-problemphyxmini0347-gasvolume}
  Pressure, volume, and absolute temperature at one equilibrium state
\end{definition}
\begin{proof}
  This is the declared model definition of Thermodynamic State.
\end{proof}
\begin{definition}[Thermodynamic Process]
  \label{def:physics:phyx-mini-0347:phyxminiproblems-problemphyxmini0347-thermodynamicprocess}
  \lean{PhyXMiniProblems.ProblemPhyXMini0347.ThermodynamicProcess}
  \uses{def:physics:phyx-mini-0347:phyxminiproblems-problemphyxmini0347-statelabel, def:physics:phyx-mini-0347:phyxminiproblems-problemphyxmini0347-processkind}
  A directed quasistatic leg. Heat transferred to the gas and work done by the gas use the convention that positive values enter the first-law identity Q = ΔU + W
\end{definition}
\begin{proof}
  This is the declared model definition of Thermodynamic Process.
\end{proof}
\begin{definition}[Ideal Gas Heat Engine Cycle]
  \label{def:physics:phyx-mini-0347:phyxminiproblems-problemphyxmini0347-idealgasheatenginecycle}
  \lean{PhyXMiniProblems.ProblemPhyXMini0347.IdealGasHeatEngineCycle}
  \uses{def:physics:phyx-mini-0347:phyxminiproblems-problemphyxmini0347-statelabel, def:physics:phyx-mini-0347:phyxminiproblems-problemphyxmini0347-thermodynamicstate, def:physics:phyx-mini-0347:phyxminiproblems-problemphyxmini0347-thermodynamicprocess}
  The ideal-gas engine cycle. The amount, gas constant, and molar heat capacity fields are explicitly named scalar readouts in mol, J/(mol K), and J/(mol K), respectively; thermalEfficiency is dimensionless
\end{definition}
\begin{proof}
  This is the declared model definition of Ideal Gas Heat Engine Cycle.
\end{proof}
\begin{definition}[Obeys First Law]
  \label{def:physics:phyx-mini-0347:phyxminiproblems-problemphyxmini0347-obeysfirstlaw}
  \lean{PhyXMiniProblems.ProblemPhyXMini0347.ObeysFirstLaw}
  \uses{def:physics:phyx-mini-0347:phyxminiproblems-problemphyxmini0347-energyinjoules, def:physics:phyx-mini-0347:phyxminiproblems-problemphyxmini0347-temperatureinkelvin, def:physics:phyx-mini-0347:phyxminiproblems-problemphyxmini0347-thermodynamicprocess, def:physics:phyx-mini-0347:phyxminiproblems-problemphyxmini0347-idealgasheatenginecycle}
  The first law for one leg, with the ideal-gas internal-energy change ΔU = n Cᵥ (T\_final - T\_initial)
\end{definition}
\begin{proof}
  This is the declared model definition of Obeys First Law.
\end{proof}
\begin{definition}[Has Physical Heat Engine Parameters]
  \label{def:physics:phyx-mini-0347:phyxminiproblems-problemphyxmini0347-hasphysicalheatengineparameters}
  \lean{PhyXMiniProblems.ProblemPhyXMini0347.HasPhysicalHeatEngineParameters}
  \uses{def:physics:phyx-mini-0347:phyxminiproblems-problemphyxmini0347-pressureinpascals, def:physics:phyx-mini-0347:phyxminiproblems-problemphyxmini0347-volumeincubicmeters, def:physics:phyx-mini-0347:phyxminiproblems-problemphyxmini0347-energyinjoules, def:physics:phyx-mini-0347:phyxminiproblems-problemphyxmini0347-temperatureinkelvin, def:physics:phyx-mini-0347:phyxminiproblems-problemphyxmini0347-statelabel, def:physics:phyx-mini-0347:phyxminiproblems-problemphyxmini0347-idealgasheatenginecycle}
  Positivity and heat-flow conditions selecting the heat-engine branch
\end{definition}
\begin{proof}
  This is the declared model definition of Has Physical Heat Engine Parameters.
\end{proof}
\begin{definition}[Satisfies Ideal Gas Heat Engine Laws]
  \label{def:physics:phyx-mini-0347:phyxminiproblems-problemphyxmini0347-satisfiesidealgasheatenginelaws}
  \lean{PhyXMiniProblems.ProblemPhyXMini0347.SatisfiesIdealGasHeatEngineLaws}
  \uses{def:physics:phyx-mini-0347:phyxminiproblems-problemphyxmini0347-pressureinpascals, def:physics:phyx-mini-0347:phyxminiproblems-problemphyxmini0347-volumeincubicmeters, def:physics:phyx-mini-0347:phyxminiproblems-problemphyxmini0347-energyinjoules, def:physics:phyx-mini-0347:phyxminiproblems-problemphyxmini0347-temperatureinkelvin, def:physics:phyx-mini-0347:phyxminiproblems-problemphyxmini0347-statelabel, def:physics:phyx-mini-0347:phyxminiproblems-problemphyxmini0347-idealgasheatenginecycle, def:physics:phyx-mini-0347:phyxminiproblems-problemphyxmini0347-obeysfirstlaw}
  Governing thermodynamic laws for this ideal-gas cycle. These fields state general physical relations; none mentions a numerical answer choice
\end{definition}
\begin{proof}
  This is the declared model definition of Satisfies Ideal Gas Heat Engine Laws.
\end{proof}
\begin{definition}[Has Heat Engine Figure Data]
  \label{def:physics:phyx-mini-0347:phyxminiproblems-problemphyxmini0347-hasheatenginefiguredata}
  \lean{PhyXMiniProblems.ProblemPhyXMini0347.HasHeatEngineFigureData}
  \uses{def:physics:phyx-mini-0347:phyxminiproblems-problemphyxmini0347-pressureinatmospheres, def:physics:phyx-mini-0347:phyxminiproblems-problemphyxmini0347-volumeincubicmeters, def:physics:phyx-mini-0347:phyxminiproblems-problemphyxmini0347-idealgasheatenginecycle}
  The textual parameters and primary-figure readouts. The equalities for the three processes also record the arrow directions a → b → c → a
\end{definition}
\begin{proof}
  This is the declared model definition of Has Heat Engine Figure Data.
\end{proof}
\begin{theorem}[pressure At A from adiabatic]
  \label{thm:physics:phyx-mini-0347:phyxminiproblems-problemphyxmini0347-pressureata-from-adiabatic}
  \lean{PhyXMiniProblems.ProblemPhyXMini0347.pressureAtA_from_adiabatic}
  \uses{def:physics:phyx-mini-0347:phyxminiproblems-problemphyxmini0347-pressureinatmospheres, def:physics:phyx-mini-0347:phyxminiproblems-problemphyxmini0347-idealgasheatenginecycle, def:physics:phyx-mini-0347:phyxminiproblems-problemphyxmini0347-hasphysicalheatengineparameters, def:physics:phyx-mini-0347:phyxminiproblems-problemphyxmini0347-satisfiesidealgasheatenginelaws, def:physics:phyx-mini-0347:phyxminiproblems-problemphyxmini0347-hasheatenginefiguredata}
  The pressure at a furnished by the adiabatic p V\textasciicircum{}γ relation (part (a) of the source problem), expressed in atmospheres. It is a derived conclusion, not a figure-data or governing-law premise
\end{theorem}
\begin{proof}
  The conclusion follows from the governing relations and figure readouts named in the statement of pressure At A from adiabatic.
\end{proof}
\begin{definition}[Efficiency Choice]
  \label{def:physics:phyx-mini-0347:phyxminiproblems-problemphyxmini0347-efficiencychoice}
  \lean{PhyXMiniProblems.ProblemPhyXMini0347.EfficiencyChoice}
  The four efficiencies displayed in the source, represented as fractions
\end{definition}
\begin{proof}
  This is the declared model definition of Efficiency Choice.
\end{proof}
\begin{definition}[efficiency Choice Fraction]
  \label{def:physics:phyx-mini-0347:phyxminiproblems-problemphyxmini0347-efficiencychoicefraction}
  \lean{PhyXMiniProblems.ProblemPhyXMini0347.efficiencyChoiceFraction}
  \uses{def:physics:phyx-mini-0347:phyxminiproblems-problemphyxmini0347-efficiencychoice}
  Numerical value of a displayed efficiency choice: 30\%, 31.9\%, 35.9\%, or 32.3\%
\end{definition}
\begin{proof}
  This is the declared model definition of efficiency Choice Fraction.
\end{proof}
\begin{definition}[Is Closest Efficiency Choice]
  \label{def:physics:phyx-mini-0347:phyxminiproblems-problemphyxmini0347-isclosestefficiencychoice}
  \lean{PhyXMiniProblems.ProblemPhyXMini0347.IsClosestEfficiencyChoice}
  \uses{def:physics:phyx-mini-0347:phyxminiproblems-problemphyxmini0347-efficiencychoice, def:physics:phyx-mini-0347:phyxminiproblems-problemphyxmini0347-efficiencychoicefraction}
  A displayed choice is closest to the calculated dimensionless efficiency
\end{definition}
\begin{proof}
  This is the declared model definition of Is Closest Efficiency Choice.
\end{proof}
% --- Archon declaration topology end ---
% --- Archon physics formalization source end ---
